\section{Introduction}
\label{sec:introduction}

Data centres consume an increasing share of global electricity, and
hyperscale operators have made public commitments to carbon-neutral
or net-zero operation within the decade. These pledges have motivated
substantial research on \emph{carbon-aware computing}: techniques
that exploit spatial and temporal variation in grid carbon intensity
to shift computation toward cleaner energy. The dominant strategy in
the literature is \emph{temporal shifting} of delay-tolerant work to
periods of low carbon intensity \citep{wiesner2021lets}; complementary
\emph{spatial shifting} routes work toward regions with cleaner grids
\citep{sukprasert2023quantifying,radovanovic2023google}. Both have delivered
20--45\% operational-carbon reductions for batch processing, machine
learning training, and similar long-running workloads
\citep{hanafy2023carbonscaler,sukprasert2023quantifying}.

A critical assumption underlies essentially this entire body of work:
the target workload is \emph{batch} in character. The canonical
scheduled job is minutes to hours long, delay-tolerant on the order
of hours, amenable to clean suspend-and-resume, and accompanied by
enough up-front information that the scheduler can reason globally
over a set of pending jobs. These assumptions are violated profoundly
by \emph{serverless} or \emph{Function-as-a-Service} (FaaS) workloads,
which now power a substantial fraction of cloud-native applications.
FaaS invocations are typically measured in tens of milliseconds to a
few seconds; they are often user-facing with sub-second latency
objectives; they cannot be suspended and resumed without paying the
full cost of a \emph{cold start}; their arrivals are bursty and
event-driven; and per-invocation behaviour is largely unpredictable
until execution.

This mismatch is consequential. Recent FaaS-targeted carbon-aware
schedulers --- \emph{GreenCourier} \citep{chadha2023greencourier},
\emph{Caribou} \citep{gsteiger2024caribou}, \emph{EcoLife}
\citep{jiang2024ecolife}, and \emph{CASA} \citep{qi2024casa} --- have
each begun to address the gap from a distinct angle: spatial routing of
single functions (GreenCourier), workflow-level geo-shifting
(Caribou), hardware-generation choice (EcoLife), and SLO/carbon-aware
autoscaling (CASA). None of them, however, formalizes the joint
spatial/temporal/warm-pool decision problem that FaaS schedulers
actually face, and the fundamental cold-start versus idle-energy
trade-off has not previously been characterized in closed form.

We argue that carbon-aware scheduling for FaaS is both feasible and
distinct enough from its batch counterpart to warrant first-class
treatment. Three observations motivate our approach. \emph{First},
FaaS workloads are heterogeneous in latency tolerance: a
customer-facing API request demands sub-second response, but many
event-driven functions tolerate seconds to minutes of deferral
without observable user impact. \emph{Second}, FaaS platforms already
perform fine-grained placement and routing decisions for
load-balancing and cold-start mitigation; carbon-aware policies can
be layered onto this existing machinery. \emph{Third}, warm-pool
management --- how many idle containers to keep alive in each region
--- is itself a carbon decision, because idle containers consume
non-trivial power and the choice of where to keep them warm interacts
directly with the carbon intensity of the host grid.

In this work, we present \GreenFaaS, a carbon-aware scheduling
framework for serverless workloads that unifies spatial routing,
temporal deferral, and warm-pool-aware container reuse under a single principled
policy. \GreenFaaS{} combines temporal deferral for asynchronous and
event-driven invocations, spatial routing for latency-tolerant
classes, and carbon-aware warm-pool management, all under an
SLA-tiered policy distinguishing interactive, deferrable, and
background functions. A key technical contribution is the explicit
modelling of the \emph{cold-start carbon trade-off}: keeping
containers warm in a low-carbon region can, in some regimes, cost
more carbon (through idle power) than cold-starting in a
higher-carbon region. We derive analytical break-even points for
this trade-off and incorporate them directly into the scheduler.

We evaluate \GreenFaaS{} in \textbf{GreenFaaS-sim}, a discrete-event
simulator whose design follows the abstractions of \code{faas-sim}
and \code{vessim} \citep{faassim,vessim}, driven by the real Azure
Functions 2021 per-invocation trace
\citep{zhang2021faster,azuretrace} and by real grid carbon traces from
Let's-Wait-Awhile \citep{wiesner2021lets,lwa} and from
ElectricityMaps over the same January 2021 window as the Azure
trace. The trace loaders also accept a schema-faithful synthetic
workload as a fallback when the real trace is not locally available
(\S\ref{sec:evaluation:methodology}). We compare against four
carbon-aware baselines drawn from the recent literature, evaluating
along five axes: workload intensity, number of available regions, SLA
class mix, carbon-intensity variability, and forecast accuracy.

Our key findings are threefold. \emph{First}, \GreenFaaS{}
reduces operational carbon by 68--83\% across topologies on fully-real
time-aligned data (Azure 2021 trace + ElectricityMaps Jan 2021
carbon), and by 64--72\% on synthetic-workload setups; on fully-real
data the \GreenFaaS{}--Spatial gap is just 0.04 percentage points.
\emph{Second}, \GreenFaaS{} preserves a \emph{do-no-harm} property:
in single-region and flat-carbon regimes where shifting offers no
real opportunity, \GreenFaaS{} matches FIFO byte-for-byte on both
synthetic and fully-real data, while a heuristic ablation (lacking
the cold-start trade-off correction) loses 5--102\% to FIFO in the
same scenarios. \emph{Third}, the canonical Wait-Awhile temporal
scheduler and a FaaS-adapted port of the
\citet{lechowicz2024online} provably-optimal one-way trading
algorithm both \emph{fail} on FaaS workloads (3--306\% over FIFO
depending on the regime), an empirically sharp demonstration that
batch-oriented carbon-aware techniques --- including those with
provable competitive guarantees --- do not transfer to FaaS without
addressing warm-pool idle-energy coupling. This robustness across
regimes --- rather than dominance in any single regime --- is the
central practical claim of the paper.

\paragraph{Contributions.} The paper makes the following contributions:
\begin{enumerate}
\item A formal model of carbon-aware FaaS scheduling as an online
  optimization problem with capacity, SLA, and warm-pool constraints,
  and an honest discussion of why classical competitive-ratio bounds
  do not apply (\S\ref{sec:problem}).
\item A closed-form \emph{cold-start carbon trade-off lemma}
  characterizing warm-pool versus cold-start carbon as a function of
  arrival rate, runtime and cold-start durations, active and idle
  power, warm-pool TTL, and the carbon-intensity ratio between
  regions (\S\ref{sec:tradeoff}). The same dimensionless idle-cost
  ratio governs the analogous temporal-shift decision via an
  idle-energy correction (\S\ref{sec:tradeoff:idle}).
\item The \GreenFaaS{} scheduler: a hybrid online algorithm combining
  temporal deferral, spatial routing, and lemma-driven warm-pool
  management under an SLA-tiered policy, with $O(R \cdot N)$
  per-invocation cost (\S\ref{sec:algorithm}).
\item An open-source simulator with trace loaders that accept the
  published Azure Functions 2019, Azure Functions 2021, and LWA
  carbon-intensity schemas without modification
  (\S\ref{sec:simulator}).
\item An empirical evaluation against four carbon-aware baselines and
  a five-axis sensitivity sweep establishing the do-no-harm property
  across topology, SLA-class mix, forecast accuracy, carbon
  variability, and workload intensity (\S\ref{sec:evaluation}).
\end{enumerate}

The remainder of this paper is organized as follows.
\S\ref{sec:related} surveys related work; \S\ref{sec:motivation}
motivates the problem; \S\ref{sec:problem} formalizes it;
\S\ref{sec:tradeoff} derives the trade-off lemma;
\S\ref{sec:algorithm} presents the algorithm; \S\ref{sec:simulator}
describes the simulator; \S\ref{sec:evaluation} reports results;
\S\ref{sec:discussion} discusses limitations; \S\ref{sec:conclusion}
concludes.
